\chapter{Conclusions and Future Work}\label{ch:conclusions}

\section{Conclusions}\label{sec:conc-conclusion}
This thesis started from a simple observation about tracking-by-detection systems: the settings that connect the detector and the tracker---the detector operating point and the score thresholds of the tracker---are usually fixed once for one detector and one kind of scene, although the scene changes within a video and the detector changes during the life of a system. The research problem was to adapt these settings online, from information that is available while the video is processed, without labels, without training, and with the detector and the tracker left unchanged (Section~\ref{sec:intro-problem}).

The published survey \cite{ismail2026realtime} provided the starting point. It reviewed the detectors, trackers, datasets, and metrics of real-time detection and tracking, showed that tracking by detection remains the dominant paradigm, and pointed to the dependence between detection accuracy and association reliability under real-time constraints. Together with the review of related work in Chapter~\ref{ch:background}, it identified the gap addressed by the thesis: existing adaptive methods either learn a policy for one detector, adapt one tracker from inside, or require labeled data.

The first controller, AC-MOT (Chapter~\ref{ch:acmot}), answered RQ1 for UAV video. It adapts the confidence threshold, the NMS threshold, and the input resolution of the detector every ten frames. A scene complexity index formed from the crowd cue, checks from four further inexpensive scene cues, and a three-state controller with recovery, hysteresis, and dwell rules choose one of three operating points; the values of the controller were selected by cross-validation on separate calibration sequences and then frozen. On seven VisDrone2019-MOT sequences it improved HOTA by 2.937 points with ByteTrack and 1.325 points with OATrack, and MOTA by 17.995 and 6.318 points, with fewer identity switches and false positives; paired bootstrap analysis supports these gains. Qualitative experiments on the tracked frames of a VisDrone video made this effect visible: far fewer false tracks, fewer identities per object, and a stable operating point in a crowded scene. On UAVDT it improved MOTA by 17.601 and 3.267 points without retuning, and it adds about 4.2 ms per frame on a Tesla T4.

The second controller, SCL (Chapter~\ref{ch:scl}), answered RQ2. It reads the score distribution of recent frames with nested Otsu splits and a whole-stream majority rule. On clean scores, it keeps the tracker's scores and, where they fit, its thresholds; on noisy scores, it chooses the candidates, rescores them, and sets the thresholds. A track-guided recovery gives single-stage trackers the continuation step they lack, and a host contract of four published settings lets one configuration serve different trackers. One frozen configuration restored two trackers from 0 to about 66 HOTA under halved scores, improved both single-stage trackers, including one not seen during design, on every MOT17 sequence, raised MOTA by up to 24.70 points with a noisy detector on held-out aerial data, cut the identity switches of two-stage trackers by 43\% to 56\% with replacement detectors, and kept well-matched trackers at their own accuracy. It reached higher HOTA than simple adaptive thresholds (by 2.2 to 8.6 points) and than the published ACT method, and it takes 2.95 ms per frame on an ordinary processor.

Both controllers kept their effect beyond the data used to design them, at a cost of a few milliseconds per frame, which answers RQ3: AC-MOT on UAVDT without retuning, and SCL on two predeclared external trackers and on the held-out UAVDT test split, with detectors for which the trackers were not tuned. Table~\ref{tab:conc-summary} summarizes the two controllers and their main results, and Table~\ref{tab:conc-rq} relates each research question to its principal evidence and main finding.

\begin{table}[tbp]
\centering
\caption{Summary of the two proposed controllers and their main results; all changes are relative to the same pipeline without the controller.}
\label{tab:conc-summary}
\begin{tabular}{>{\raggedright\arraybackslash}p{2.7cm}>{\raggedright\arraybackslash}p{5.6cm}>{\raggedright\arraybackslash}p{5.6cm}}
\toprule
Aspect & AC-MOT & SCL\\
\midrule
Question & RQ1: adapt the detector operating point to the scene & RQ2: restore the compatibility of score scale and tracker thresholds\\
Acts on & detector operating point & tracker input and thresholds\\
Evaluated on & VisDrone2019-MOT and UAVDT with ByteTrack and OATrack & MOT17, KITTI, and UAVDT with seven published trackers\\
Main results & VisDrone: HOTA $+2.937$ and $+1.325$, MOTA $+17.995$ and $+6.318$; UAVDT: MOTA $+17.601$ and $+3.267$ & halved scores: 0 to 66.17 and 66.72 HOTA; KITTI: MOTA $+5.98$ and $+7.86$; UAVDT: MOTA $+24.70$\\
Cost & 4.26 and 4.20 ms per frame (Tesla T4) & 2.95 ms per frame (CPU)\\
\bottomrule
\end{tabular}
\end{table}

\begin{table}[tbp]
\centering
\caption{Research questions of the thesis, the principal experiments that answer them, and the main findings.}
\label{tab:conc-rq}
\begin{tabular}{>{\raggedright\arraybackslash}p{2.9cm}>{\raggedright\arraybackslash}p{5.3cm}>{\raggedright\arraybackslash}p{5.7cm}}
\toprule
Research question & Principal evidence & Main finding\\
\midrule
RQ1: online control of the detector operating point & AC-MOT on seven VisDrone2019-MOT sequences with ByteTrack and OATrack, and frozen transfer to UAVDT (Sections~\ref{sec:ac-res-visdrone}--\ref{sec:ac-res-uavdt}) & over the fixed operating point: HOTA $+2.937$ and $+1.325$ and MOTA $+17.995$ and $+6.318$ on VisDrone; MOTA $+17.601$ and $+3.267$ on UAVDT without retuning\\
RQ2: label-free compatibility of scores and tracker thresholds & SCL under halved scores, with replacement detectors, and on MOT17, KITTI, and UAVDT, compared with adaptive thresholds and ACT (Sections~\ref{sec:scl-results} and~\ref{sec:scl-compare}) & two trackers restored from 0 to about 66 HOTA; well-matched trackers kept at their own accuracy; higher HOTA than simple adaptive thresholds (by 2.2 to 8.6 points) and than ACT\\
RQ3: effect on unseen data, trackers, and detectors, and cost & AC-MOT on UAVDT without retuning; SCL on two predeclared external trackers and the held-out UAVDT test split; runtime (Sections~\ref{sec:ac-res-uavdt}, \ref{sec:ac-res-runtime}, \ref{sec:scl-res-external}, \ref{sec:scl-res-uavdt}, and~\ref{sec:scl-runtime}) & effects kept on held-out data and on trackers not used in design; 4.2 ms per frame for AC-MOT (Tesla T4) and 2.95 ms per frame for SCL (CPU)\\
\bottomrule
\end{tabular}
\end{table}

Taken together, the two studies advance the understanding of robust tracking in two ways. First, they show that the operating point of the detector and the score scale that the tracker receives are parts of the tracking system that deserve the same attention as the association method: with the detector and the tracker unchanged, each of them moved MOTA or HOTA by many points. Second, they show that both can be controlled from outside the detector and the tracker, with rules whose values were fixed before the final evaluation: AC-MOT selects the detector settings that suit the current aerial scene, and SCL makes the scores of a given detector usable by different trackers and leaves a well-matched pair unchanged. The two controllers are distinct methods that answer different questions: AC-MOT acts before the detector and does not change the thresholds of the tracker, and SCL acts between the detector and the tracker and does not choose the detector operating point. Each was developed, frozen, and evaluated on its own, AC-MOT in Chapter~\ref{ch:acmot} and SCL in Chapter~\ref{ch:scl}.

\emph{Practical implications.} Neither controller needs training, labels, or access to the inside of the detector or the tracker. AC-MOT changes three arguments of the detector call, and SCL sets the published thresholds and the input of the tracker for each frame. A system integrator can therefore keep a tuned tracker when the scenes change or when the detector is replaced or retrained, and the label-free mode of SCL can be logged as a running check that a new detector still fits its tracker.

\section{Research Limitations}\label{sec:conc-limitations}
The results of Chapters~\ref{ch:acmot} and~\ref{ch:scl} hold for the conditions under which they were obtained. This section states five limitations that bound their interpretation: two concern the signals from which the controllers decide, two concern the range of detectors and trackers that the evidence covers, and one concerns deployment hardware. Each limitation is stated together with what the experiments did show, so that the demonstrated behavior is not confused with untested conditions.

\subsection{Dependence on Reliable Scene Cues}\label{sec:lim-cues}
AC-MOT decides from inexpensive scene cues: the crowd cue and the tiny-object ratio, both computed from the boxes that the tracker reported in the previous frame, and the edge density, brightness, and sharpness of the current image (Section~\ref{sec:ac-cues}). Its state is therefore only as good as these cues. Because the crowd cue counts tracked boxes, it measures the scene through the output of the tracker: a tracker that reports many false tracks makes a scene look more crowded than it is, and a tracker that misses most objects, for example at night or under heavy blur, makes it look sparser. The scene label uses fixed thresholds on brightness and sharpness, and these thresholds were chosen for the cameras and conditions of the development data. In such cases the controller may select an operating point that does not suit the actual scene. The design reduces the effect of short disturbances, since the SCI averages seven analysis steps, the hysteresis and dwell rules prevent rapid switching, and the recovery probe reacts to a sudden collapse of the object count (Section~\ref{sec:ac-states}). On the evaluated sequences the controller used all three states and changed state rarely (Section~\ref{sec:ac-res-behavior}). The experiments did not, however, test sequences in which the cues are deliberately corrupted, and the behavior of the controller under strongly misleading cues remains an open question rather than a demonstrated failure or success.

\subsection{Sensitivity to Detection-Score Distributions}\label{sec:lim-scores}
SCL reads the distribution of recent detection scores and assumes that this distribution carries information about the objects: that the order of the scores separates objects from clutter well enough for the nested Otsu splits to form meaningful background, uncertain, and confident groups (Sections~\ref{sec:scl-calib} and~\ref{sec:scl-regimes}). This assumption held for the detectors, score transforms, and datasets of Chapter~\ref{ch:scl}, including halved, cubed, and temperature-scaled scores. The splits are always computed, however, even when the scores are poorly separated, for example when the scores of objects and clutter overlap strongly, when a detector emits only a narrow range of scores, or when a frame window contains very few candidates. In such cases the groups may not correspond to objects and clutter, and the confident share $\rho$ and the clean or noisy decision become less informative. The whole-stream median keeps the mode stable against short disturbances, and the layer falls back to the tracker's own thresholds when no valid split exists (the cold mode, Algorithm~\ref{alg:scl-layer}), but a stream whose score distribution changes its character permanently in the middle of a video would be described by a median that adapts only slowly. These are assumptions of the score-based design; the experiments did not include detectors whose scores are uninformative, and no failure of this kind was observed in the evaluated settings.

\subsection{Limited Detector Coverage of AC-MOT}\label{sec:lim-detectors}
The frozen AC-MOT controller was evaluated with one detector, YOLO11m trained on VisDrone2019-DET, on two aerial datasets, VisDrone2019-MOT and UAVDT, and with two trackers, ByteTrack and OATrack (Section~\ref{sec:ac-setup}). The earlier development stages used YOLOv8n with ByteTrack, the YOLOX-X detector of U2MOT, and SparseTrack on MOT17 (Section~\ref{sec:ac-dev}), but with earlier versions of the controller and other protocols, so they do not extend the evidence for the frozen controller. Its action map is also bound to the operating points for which detector outputs were cached: input resolutions of 1088, 1280, and 1536 pixels, NMS thresholds of 0.45, 0.60, and 0.70, and score thresholds of 0.10, 0.25, and 0.40. The improvements reported in Chapter~\ref{ch:acmot} therefore establish the benefit of the frozen controller for this detector and these operating points. They do not by themselves establish the same benefit for detectors of other families, such as two-stage or transformer-based detectors, whose score distributions and sensitivity to input resolution differ, or for detectors trained on other data. For a new detector, the controller would at least need to be checked, and possibly recalibrated with the same cross-validation procedure.

\subsection{The Host Contract of SCL}\label{sec:lim-contract}
SCL serves a tracker through a host contract: the tracker must publish four values, the matching threshold $a$, the birth threshold $b$, the lowest score $\ell$ that it uses, and the matching tolerance $m_0$, and it must accept, for every frame, a list of candidate boxes with scores together with new values of these thresholds and of the tolerance (Section~\ref{sec:scl-contract}). The architectural requirements behind this contract are that the tracker separates detection from association, reads the detector scores through explicit thresholds, and matches tracks and boxes with an overlap-based tolerance that can be set from outside. Trackers that meet these requirements were served by one fixed configuration in this thesis: seven published tracking-by-detection trackers, including two external trackers chosen before the run. Other parts of the integration are implementation choices: each tracker family needed a small adapter that writes the values of the layer into its existing parameters before each update, so a new tracker requires a new adapter even when it satisfies the contract. Trackers that do not expose such control points, for example joint detection-and-tracking and transformer-based trackers whose thresholds are internal to a network (Section~\ref{sec:bg-paradigms}), or trackers that associate mainly by appearance without score thresholds, cannot be served directly. SCL is therefore applicable to trackers with an explicit score-threshold interface, not to every tracker.

\subsection{Incomplete Embedded-Hardware Validation}\label{sec:lim-hardware}
The computational cost of both controllers was measured, but not on embedded hardware. For AC-MOT, the time spent in the controller was measured on a Tesla T4 GPU over 200 frames (Section~\ref{sec:ac-res-runtime}); for SCL, the time of the layer and the end-to-end time of a YOLOv8n and ByteTrack pipeline were measured on a 4-vCPU server processor without a GPU (Section~\ref{sec:scl-runtime}). These measurements show that the controllers themselves are inexpensive compared with the detector, but they are not an end-to-end validation on the resource-constrained processors carried by UAVs. On such platforms, the detector at each operating point determines most of the latency and the throughput, and AC-MOT changes the input resolution, so the end-to-end frame rate depends on the hardware. Memory use, power consumption, thermal limits during long flights, and the effect of hardware-specific optimizations such as reduced numerical precision were not measured, and the sustained performance of the complete pipeline over long operation was not evaluated. The real-time suitability shown in this thesis therefore refers to the measured cost of the controllers, not to a demonstrated deployment on embedded hardware.

\section{Future Work}\label{sec:conc-future}
The limitations above point to the following directions, which build on the training-free and modular design of the two controllers. None of these studies has been carried out in this thesis.
\begin{enumerate}
\item \emph{Robust scene cues and controller reliability} (Section~\ref{sec:lim-cues}). The scene cues of AC-MOT can be made less dependent on the output of the tracker, for example by combining the count of tracked boxes with cues that the detector provides directly, and the reliability of the state decisions can be studied on sequences with abrupt changes of illumination, blur, and crowd density, by recording the state changes per 100 frames together with HOTA, MOTA, and IDF1.
\item \emph{Score-distribution robustness of SCL} (Section~\ref{sec:lim-scores}). A label-free measure of how well the score distribution is separated could tell the layer when its splits are reliable and when it should keep the tracker's own thresholds, and alternative label-free partitions of the scores, such as mixture models, can be compared with the nested Otsu splits on detectors with compressed or poorly separated score distributions.
\item \emph{More detectors and datasets for AC-MOT} (Section~\ref{sec:lim-detectors}). Because AC-MOT needs no training and changes only three arguments of the detector call, the same calibration procedure can be applied to other detector families, including two-stage and transformer-based detectors, and to further aerial datasets, and its action map can include further operating points. Each new detector is to be evaluated in two separate conditions, the frozen controller transferred without retuning and the controller recalibrated with the same cross-validation procedure, by HOTA, MOTA, and IDF1 against the detector at its fixed operating point.
\item \emph{Portability of SCL to other tracker interfaces} (Section~\ref{sec:lim-contract}). The host contract can be extended to trackers that learn detection and association jointly or that rely mainly on appearance, by identifying the score thresholds or confidence gates that such trackers use internally, and a generic adapter interface would reduce the work of adding a new tracker. As in Chapter~\ref{ch:scl}, each extension is to be judged by HOTA and IDS against the same tracker alone under score shifts and detector changes.
\item \emph{Embedded-hardware evaluation} (Section~\ref{sec:lim-hardware}). The complete pipeline, with each controller, is to be evaluated on embedded processors of the kind carried on UAVs, such as an NVIDIA Jetson module, with a target of 30 FPS, that is, at most 33.3 ms per frame for the detector, the controller, and the tracker together, and with latency, throughput, memory use, and power consumption measured over long runs.
\item \emph{Deployment optimizations.} On such hardware, the cost of the detector at each operating point of AC-MOT, the use of reduced numerical precision, and the scheduling of scene analysis and tracking can be optimized for latency, memory, throughput, and energy, and the effect of each optimization on the decisions of the controllers and on tracking accuracy is to be measured.
\item \emph{Monitoring of detector updates.} The label-free mode and the confident share of SCL can serve as a running check when the detector of a deployed system is updated, retrained, or compressed; the objective is to log the mode and $\rho$ before and after each update and to relate them to the tracking accuracy measured on labeled test sequences.
\item \emph{A combined pipeline.} Since the two controllers act on different parts of the pipeline, they can be run together, with AC-MOT setting the detector operating point and SCL adapting the tracker input to the resulting scores, to evaluate whether their improvements are complementary. The planned evaluation compares four configurations on the same sequences---the original pipeline, AC-MOT alone, SCL alone, and both controllers together---by tracking accuracy (HOTA and MOTA), identity preservation (IDF1 and IDS), the stability of the controller states and of the SCL modes (changes per 100 frames), and the processing time per frame.
\end{enumerate}
